% 第 11 题四棱柱三视图：按固定源 PDF 的视图排布与尺寸线位置重绘（数字只保留原图印出的 1、1、1、2）。
% 正（主）视图为 2:1 矩形、内部两条竖直虚线在宽度的 1/4 与 3/4 处（梯形两斜棱的投影），原图无尺寸标注。
% 俯视图整块较视图标签下移 0.5cm：标签（锚 y=-0.32，字高约 0.37）与上底尺寸线（y=-1.12）留出 0.43cm 白隙。
\begin{tikzpicture}[every node/.style={font=\small},scale=0.9,line cap=round,line join=round]
% 正（主）视图：矩形内两条竖直虚线，无尺寸标注
\draw (0,0) rectangle (2.0,1.0);
\draw[dashed] (0.5,0)--(0.5,1.0);
\draw[dashed] (1.5,0)--(1.5,1.0);
\node[below] at (1.0,-0.32) {正（主）视图};
% 侧（左）视图：正方形；右侧竖直尺寸线（两端端线），数字 1 在旁
\draw (3.3,0) rectangle (4.3,1.0);
\draw[<->] (4.68,0)--(4.68,1.0);
\draw (4.55,0)--(4.81,0); \draw (4.55,1.0)--(4.81,1.0);
\node[right] at (4.85,0.5) {\(1\)};
\node[below] at (3.8,-0.32) {侧（左）视图};
% 俯视图：等腰梯形（上底 1、下底 2、高 1）；上底上方数字 1 断开的水平尺寸线，
% 右侧竖直尺寸线旁标 1，下底下方数字 2 断开的水平尺寸线
\draw (0,-2.4)--(2.0,-2.4)--(1.5,-1.4)--(0.5,-1.4)--cycle;
\draw[<->] (0.5,-1.12)--(1.5,-1.12);
\node[fill=white,inner sep=1.2pt] at (1.0,-1.12) {\(1\)};
\draw[<->] (2.42,-1.4)--(2.42,-2.4);
\draw (2.29,-1.4)--(2.55,-1.4); \draw (2.29,-2.4)--(2.55,-2.4);
\node[right] at (2.58,-1.9) {\(1\)};
\draw[<->] (0,-2.68)--(2.0,-2.68);
\node[fill=white,inner sep=1.2pt] at (1.0,-2.68) {\(2\)};
\node[below] at (1.0,-3.05) {俯视图};
\end{tikzpicture}
